%% ============================================================
%% Hauptteil
%% Der Hauptteil kann in mehrere Kapitel aufgeteilt werden.
%% Diese Datei enthaelt exemplarisch drei Kapitel:
%%   - Theoretische Grundlagen
%%   - Analyse / Untersuchung
%%   - Ergebnisse
%%
%% Bei Bedarf koennen die Kapitel auch in eigene Dateien ausgelagert
%% und einzeln per \input{} eingebunden werden.
%% ============================================================


% ============================================================
\chapter{Theoretische Grundlagen}
\label{chap:theoretische-grundlagen}
% ============================================================

% Hinweis: Dieses Kapitel stellt den wissenschaftlichen Hintergrund dar.
% Hier werden relevante Theorien, Begriffe und der Stand der Forschung
% erlaeutert. Alle verwendeten Quellen muessen korrekt zitiert werden.

[Einfuehrung in den theoretischen Rahmen der Arbeit. Welche Begriffe
und Konzepte sind zentral?]


\section{Begriffsbestimmung}
\label{sec:begriffe}

[Definition und Abgrenzung der wichtigsten Begriffe, die in der Arbeit
verwendet werden.]

% Beispiel fuer ein indirektes Zitat (Paraphrase) im APA-Stil:
% Laut \textcite{nachname2023} ist [Inhalt der Paraphrase].
% Oder: [Inhalt der Paraphrase] \parencite{nachname2023}.


\section{Stand der Forschung}
\label{sec:forschungsstand}

[Darstellung des aktuellen Forschungsstandes zum Thema. Welche
Arbeiten und Erkenntnisse gibt es bereits? Wo bestehen Forschungsluecken?]

% Beispiel fuer ein direktes Zitat im APA-Stil:
% \textcite[S.~42]{nachname2023} stellt fest: \glqq [Zitat]\grqq.


\section{Theoretischer Rahmen}
\label{sec:theorie}

[Vorstellung der theoretischen Grundlage, auf der die Arbeit aufbaut.
Welche Modelle oder Theorien werden herangezogen?]


% ============================================================
\chapter{Analyse}
\label{chap:analyse}
% ============================================================

% Hinweis: In diesem Kapitel wird die eigentliche Untersuchung
% durchgefuehrt. Hier werden die Methoden angewendet und die
% gesammelten Daten oder Quellen ausgewertet.

[Einfuehrung in das Analysekapitel. Was wird untersucht und wie?]


\section{Materialauswahl und Vorgehensweise}
\label{sec:material}

[Beschreibung der verwendeten Materialien (Quellen, Daten, Texte)
und der konkreten Vorgehensweise bei der Analyse.]


\section{Durchfuehrung der Analyse}
\label{sec:durchfuehrung}

[Darstellung der eigentlichen Analyse. Hier werden Argumente
entwickelt, Quellen interpretiert oder Daten ausgewertet.]

% Beispiel fuer eine Abbildung:
% \begin{figure}[htbp]
%   \centering
%   \includegraphics[width=0.8\textwidth]{beispiel-grafik}
%   \caption{Beschreibung der Abbildung}
%   \label{fig:beispiel}
% \end{figure}

% Beispiel fuer eine Tabelle:
% \begin{table}[htbp]
%   \centering
%   \caption{Beschreibung der Tabelle}
%   \label{tab:beispiel}
%   \begin{tabular}{lcc}
%     \toprule
%     \textbf{Kategorie} & \textbf{Wert A} & \textbf{Wert B} \\
%     \midrule
%     Zeile 1 & 42 & 58 \\
%     Zeile 2 & 37 & 63 \\
%     \bottomrule
%   \end{tabular}
% \end{table}


% ============================================================
\chapter{Ergebnisse}
\label{chap:ergebnisse}
% ============================================================

% Hinweis: Dieses Kapitel praesentiert die Ergebnisse der Analyse.
% Die Ergebnisse werden strukturiert dargestellt und interpretiert.

[Einfuehrung: Was sind die zentralen Ergebnisse der Untersuchung?]


\section{Darstellung der Ergebnisse}
\label{sec:darstellung}

[Systematische Praesentation der Ergebnisse. Jedes Teilergebnis
wird klar benannt und mit Belegen untermauert.]


\section{Interpretation und Diskussion}
\label{sec:diskussion}

[Einordnung der Ergebnisse in den theoretischen Kontext.
Was bedeuten die Ergebnisse? Wie verhalten sie sich zum
Stand der Forschung? Gibt es ueberraschende Befunde?]

[Kritische Reflexion der eigenen Vorgehensweise und moegliche
Einschraenkungen der Ergebnisse.]
